\subsubsection*{\underline{Lagrangian Relaxation:} Question 5 }

Let $I = \{1,\dots,15\}$ be the set of initiatives and $J = \{1,2,3\}$ the set
of teams. For initiative $i \in I$, $b_i$ is its strategic benefit and $a_i$
its resource consumption (Table~1). Team $j \in J$ has capacity $C_j$, with
$C_1 = 28$, $C_2 = 32$ and $C_3 = 35$. Let
\[
P = \{(2,1),(3,1),(4,1),(7,5),(8,5),(9,6),(10,6),(11,7),(12,8),(15,10)\}
\]
be the set of precedence pairs, where $(i,k) \in P$ means that initiative $i$
requires initiative $k$.

Let $x_{ij} \in \{0,1\}$ equal 1 if initiative $i$ is selected and assigned to team $j$, and $y_i \in \{0,1\}$ equal 1 if initiative $i$ is selected. So the model is:
\begin{alignat}{2}
\max \quad & \sum_{i \in I} b_i\, y_i && \label{eq:q5obj}\\
\text{s.t.} \quad
& \sum_{j \in J} x_{ij} = y_i && \qquad \forall i \in I \label{eq:q5assign}\\
& \sum_{i \in I} a_i\, x_{ij} \le C_j && \qquad \forall j \in J \label{eq:q5cap}\\
& y_i \le y_k && \qquad \forall (i,k) \in P \label{eq:q5prec}\\
& y_{11} + y_{12} \le 1 && \label{eq:q5excl}\\
& y_2 + y_3 + y_4 + y_7 + y_8 = 3 && \label{eq:q5card}\\
& x_{ij},\, y_i \in \{0,1\} && \qquad \forall i \in I,\ j \in J \label{eq:q5dom}
\end{alignat}

\eqref{eq:q5obj} maximises the total strategic benefit. \eqref{eq:q5assign}
assigns every selected initiative to exactly one team and unselected
initiatives to none. \eqref{eq:q5cap} enforces the team capacities.
\eqref{eq:q5prec} ensures that an initiative is only selected if its
prerequisite is selected. \eqref{eq:q5excl} is the exclusivity requirement, and
\eqref{eq:q5card} selects exactly three initiatives from $\{2,3,4,7,8\}$.

Without the side constraints \eqref{eq:q5prec}--\eqref{eq:q5card}, the model is
exactly the \emph{Multiple Knapsack Problem} (MKP): each initiative is an item
with profit $b_i$ and weight $a_i$, each team is a knapsack with capacity
$C_j$, and each item is placed in at most one knapsack. The total resource
consumption of all initiatives is $148$, while the total capacity is only
$28 + 32 + 35 = 95$, so the capacity constraints are binding and a selection
has to be made. Unlike the Generalized Assignment Problem, items need not be
assigned, and profits and weights do not depend on the team. The problem is
therefore a Multiple Knapsack Problem with additional precedence, conflict and
cardinality side constraints.



\subsubsection*{Question 6}
We relax the side constraints \eqref{eq:q5prec}--\eqref{eq:q5card} and keep
the assignment and capacity constraints \eqref{eq:q5assign}--\eqref{eq:q5cap}.
The multipliers are
\begin{itemize}
  \item $\lambda_{ik} \ge 0$ for the precedence constraints
        $y_i - y_k \le 0$, $(i,k) \in P$,
  \item $\nu \ge 0$ for the exclusivity constraint $y_{11} + y_{12} \le 1$,
  \item $\mu \in \mathbb{R}$ (free, since it is an equality) for the
        cardinality constraint $\sum_{i \in G} y_i = 3$, where
        $G = \{2,3,4,7,8\}$.
\end{itemize}

The relaxed model then becomes:
\begin{alignat}{2}
Z_{LR}(\lambda,\nu,\mu) = \max \quad
& \sum_{i \in I} b_i y_i
  - \sum_{(i,k) \in P} \lambda_{ik}(y_i - y_k)
  - \nu\,(y_{11} + y_{12} - 1)
  + \mu \Big(3 - \sum_{i \in G} y_i\Big) && \label{eq:q6obj}\\
\text{s.t.} \quad
& \sum_{j \in J} x_{ij} = y_i &&  \forall i \in I \notag\\
& \sum_{i \in I} a_i x_{ij} \le C_j &&  \forall j \in J \notag\\
& x_{ij},\, y_i \in \{0,1\} &&  \forall i \in I,\ j \in J \notag
\end{alignat}

Collecting terms, the objective can be written as
\[
Z_{LR}(\lambda,\nu,\mu) = \nu + 3\mu + \max \sum_{i \in I} \tilde b_i\, y_i,
\]
with modified benefits
\[
\tilde b_i = b_i
- \sum_{k:(i,k) \in P} \lambda_{ik}
+ \sum_{k:(k,i) \in P} \lambda_{ki}
- \nu\, [\, i \in \{11,12\} \,]
- \mu\, [\, i \in G \,],
\]
where $[\cdot]$ equals 1 if the condition holds and 0 otherwise.

For any $\lambda \ge 0$, $\nu \ge 0$ and $\mu \in \mathbb{R}$, every feasible
solution of the original problem is feasible in the relaxation and the penalty
terms are non-negative, so $Z_{LR}(\lambda,\nu,\mu) \ge Z_{IP}$.
The remaining subproblem is a Multiple Knapsack Problem with profits
$\tilde b_i$; initiatives with $\tilde b_i \le 0$ are simply not selected.
This subproblem does not have the integrality property, since its LP relaxation
can have fractional solutions due to the capacity constraints. Therefore the
Lagrangian dual bound can be strictly tighter than the LP relaxation bound,
$Z_{IP} \le Z_{LD} \le Z_{LP}$, which is the purpose of the investigation.


\subsubsection*{Question 7}

The Lagrangian dual problem finds the multipliers that give the tightest
(smallest) upper bound from the relaxation in Question 6:
\begin{equation}
Z_{LD} = \min_{\lambda,\, \nu,\, \mu} \; Z_{LR}(\lambda,\nu,\mu)
\quad \text{s.t.} \quad
\lambda_{ik} \ge 0 \ \ \forall (i,k) \in P, \qquad
\nu \ge 0, \qquad
\mu \in \mathbb{R}.
\label{eq:q7dual}
\end{equation}
Written out in full, this is
\[
Z_{LD} = \min_{\lambda \ge 0,\, \nu \ge 0,\, \mu \in \mathbb{R}}
\left\{ \nu + 3\mu +
\max_{(x,y) \in X} \sum_{i \in I} \tilde b_i(\lambda,\nu,\mu)\, y_i \right\},
\]
where $X$ is the set of solutions $(x,y)$ that satisfy the assignment and
capacity constraints \eqref{eq:q5assign}--\eqref{eq:q5cap} and are binary, and
$\tilde b_i(\lambda,\nu,\mu)$ are the modified benefits from Question 6.

$Z_{LR}$ is the maximum of a finite number of functions that are linear in
$(\lambda,\nu,\mu)$, one for each $(x,y) \in X$. It is therefore convex and
piecewise linear, but not differentiable everywhere, which is why we use a
subgradient method to solve \eqref{eq:q7dual}. By weak duality,
$Z_{IP} \le Z_{LD}$, and since the subproblem does not have the integrality
property, $Z_{LD} \le Z_{LP}$, possibly with strict inequality.




\section*{Question 8}

The LP relaxation of the model in Question 5 and the Lagrangian dual from
Question 7 were solved in Python with SciPy/HiGHS. Each Lagrangian subproblem, a Multiple Knapsack
Problem with modified benefits $\tilde b_i$, was solved to optimality as a
MIP. The multipliers were initialised at zero and updated with Polyak steps
\[
t_k = \theta_k \frac{Z_{LR}(\lambda^k,\nu^k,\mu^k) - \underline{Z}}{\lVert g^k \rVert^2},
\]
where $g^k$ is the violation of the relaxed constraints in the subproblem
solution, $\underline{Z} = 244$ is the best known feasible value, and
$\theta_0 = 2$ is halved after 10 iterations without improvement.
Multipliers of inequality constraints are projected onto $\mathbb{R}_+$.

\begin{center}
\begin{tabular}{@{}lr@{}}
\toprule
Bound & Value \\
\midrule
Optimal IP value $Z_{IP}$ & $244.00$ \\
Lagrangian dual bound $Z_{LD}$ (subgradient) & $254.73$ \\
LP relaxation bound $Z_{LP}$ & $254.89$ \\
\bottomrule
\end{tabular}
\end{center}

The best upper bound obtained is $Z_{LD} \approx 254.73$. The exact value,
verified by column generation, is $254.728$. It is slightly tighter than the
LP bound, as predicted by $Z_{IP} \le Z_{LD} \le Z_{LP}$, but the improvement
is only about $0.16$. Since all benefits are integer, both bounds give the same
integer bound of $254$, and the gap to $Z_{IP} = 244$ remains about $10.7$.





\section*{Question 9}

The Lagrangian relaxation does \emph{not} terminate with a feasible solution.
The subproblem solution in the final subgradient iteration
($Z_{LR} = 254.73$) is shown in Table~\ref{tab:q9lr}.

\begin{table}[h]
\centering
\begin{tabular}{@{}lll@{}}
\toprule
Team (capacity) & Initiatives & Load \\
\midrule
1 (28) & 7, 13, 14 & 28 \\
2 (32) & 5, 10, 12 & 32 \\
3 (35) & 2, 3, 6, 8 & 35 \\
\bottomrule
\end{tabular}
\caption{Subproblem solution in the final subgradient iteration
(benefit 251).}
\label{tab:q9lr}
\end{table}

It violates the precedence constraints for initiatives 2 and 3 (initiative 1
is not selected), and the cardinality constraint, since four initiatives
(2, 3, 7, 8) are selected from $G$.

\paragraph{Lagrangian heuristic.}
Starting from the subproblem solution, including its team assignment:
\begin{enumerate}
  \item \textbf{Precedence:} drop every initiative whose prerequisite is not
        selected, together with any initiatives that depend on it.
  \item \textbf{Exclusivity:} if both 11 and 12 are selected, drop the one
        with the lower ratio $b_i/a_i$.
  \item \textbf{Cardinality:} if more than three initiatives from $G$ are
        selected, drop those with the lowest benefit. If fewer, add the
        initiative from $G$ with the highest $b_i/a_i$ that fits, together
        with any missing prerequisites, placing each in the team with the
        most free capacity. If none fits, drop the selected non-$G$
        initiative with the lowest $b_i/a_i$ that is not a prerequisite of
        another selected initiative, and try again.
  \item \textbf{Greedy fill:} add the remaining non-$G$ initiatives in
        decreasing order of $b_i/a_i$ whenever they fit in a team and their
        prerequisites and exclusivity allow it.
\end{enumerate}
Steps 1--3 only remove initiatives or add initiatives that fit, so the
capacity constraints remain satisfied throughout and the result is feasible.

\paragraph{Demonstration.}
\begin{enumerate}
  \item Initiatives 2 and 3 are dropped, since initiative 1 is missing.
        Team 3 now holds $\{6, 8\}$ with load 20, and $G$ contains only
        $\{7, 8\}$.
  \item Initiative 11 is not selected, so nothing changes.
  \item One more initiative from $G$ is needed. The best candidate is 2
        ($b/a = 3.0$), which requires initiative 1, i.e.\ $22$ units, but
        only 15 units are free. Initiative 10 has the lowest ratio among
        removable non-$G$ initiatives ($25/8$), so it is dropped, freeing
        8 units in team 2. Initiative 1 is then placed in team 3 and
        initiative 2 in team 2.
  \item Only one unit of capacity remains, so no initiative can be added.
\end{enumerate}

\begin{table}[H]
\centering
\begin{tabular}{@{}lll@{}}
\toprule
Team (capacity) & Initiatives & Load \\
\midrule
1 (28) & 7, 13, 14 & 28 \\
2 (32) & 2, 5, 12 & 32 \\
3 (35) & 1, 6, 8 & 34 \\
\bottomrule
\end{tabular}
\caption{Feasible solution from the Lagrangian heuristic (benefit 218).}
\label{tab:q9heur}
\end{table}

The resulting solution in Table~\ref{tab:q9heur} satisfies all constraints
and has benefit 218. Applying the heuristic in every subgradient iteration
gives a best feasible value of 220, so
\[
220 \le Z_{IP} \le 254.
\]
The optimal value $Z_{IP} = 244$ lies within this interval. The heuristic is simple and myopic: the relaxation fills all three teams completely, so repairing the precedence violations forces the removal of valuable initiatives. A local search with swaps of up to two initiatives in and out did not improve the solution.